\chapter*{A Door into Chance}
\addcontentsline{toc}{chapter}{A Door into Chance}

We have spent a whole volume asking which changes preserve structure. Now suppose the change is random.

Take the group $\Z/2\Z=\{0,1\}$. Start at $0$. At each step, stay where you are with probability $1-p$ and add $1$ with probability $p$, where $0\le p\le1$. After $t$ steps, where are you likely to be?

Let $u_t$ be the probability of being at $0$ after $t$ steps. To return to $0$, you can stay at $0$ or move there from $1$. These are disjoint possibilities, so
\[
u_{t+1}=(1-p)u_t+p(1-u_t)
       =p+(1-2p)u_t,\qquad u_0=1.
\]
Subtract $\frac12$ from both sides:
\[
u_{t+1}-\frac12=(1-2p)\Bigl(u_t-\frac12\Bigr).
\]
The same factor appears at every step. Induction now gives
\[
\boxed{\displaystyle
\Pr(X_t=0)=u_t=\frac{1+(1-2p)^t}{2}.}
\]
Check the edges. If $p=0$, the position never changes. If $p=1$, it alternates between $0$ and $1$. If $0<p<1$, the factor $(1-2p)^t$ tends to $0$ and the distribution tends to one half at each position. In particular, $p=\frac12$ gives the uniform distribution after just one step.

The two numbers $1$ and $1-2p$ are also the eigenvalues of the transition matrix
\[
\begin{pmatrix}1-p&p\\p&1-p\end{pmatrix}.
\]
The constant vector records total probability. The other eigenvector, $(1,-1)$, measures how far we are from uniform. That is the simplest example of Fourier analysis on a group turning random motion into a power of a number.

A fair coin, a random walk, and a shuffled deck all ask the next question: when does repeated randomness wash away the starting point, and how fast? Volume~4, \emph{Chance: The Logic of Uncertainty}, begins with counting outcomes and builds the tools for answering that question.
